% !TEX root = main.tex
\documentclass[main.tex]{subfiles}
\begin{document}

% DRAFT (Section 5: constrained / penalty / discussion).
% Previous versions: flow_v1.tex (implicit / semi-implicit), flow_v2.tex, flow_v3.tex, flow_v4.tex (holonomic / penalty / projection / prior work), flow_v5.tex (constrained / implicit / penalty / discussion).

\section{Power Particles Revisited}
\label{sec:flow}

The Power Particles method \cite{deGoes2015PowerParticles} introduced a discretization of incompressible flows as moving power diagrams with prescribed volumes.
In this section, we revisit this discretization presenting a formulation based on a variational principle, which we obtain by building on the constrained power diagrams of \S\ref{sec:contributions}.
We first model the flow as a constrained Lagrangian system (\S\ref{sec:flow-dynamics}), together with its variational integrator (\S\ref{sec:flow-constrained}).
We then propose a relaxation of the constraint that keeps the variational structure while making each time step an unconstrained optimization (\S\ref{sec:flow-penalty}).
Finally, we discuss these results and relate them to prior work (\S\ref{sec:flow-discussion}).

\subsection{Constrained Dynamics}
\label{sec:flow-dynamics}

Following Power Particles \cite{deGoes2015PowerParticles}, we discretize the fluid into $n$ parcels, each one a cell $\cell_i$ of a power diagram with prescribed volume $\pvol_i$.
Unlike Power Particles, which attaches fluid particles to the sites of the power diagram, we represent each fluid parcel by a particle $\qq_i$ located \emph{exactly} at the centroid of its cell.
% The flow thus tracks the centroids of the cells rather than the cells themselves.
For simplicity, we assume a uniform density of one so that each particle has mass $\pvol_i$.
To this end, we introduce a potential $\cpot(\qq)$ that quantifies the cost of this spatial discretization by measuring how far the fluid within each cell lies from its particle, subject to volume constraints, which make the flow preserve volume, and centroid constraints, which keep the particles at the cell centroids.
% We thus define this quantization potential as
% \begin{subequations}
%   \label{eq:flow-quantization-cost}
%   \begin{alignat}{2}
%     & \cpot(\qq) = \min_{\cell} \quad && \sum\nolimits_i \int_{\cell_i} \norm{\xx - \qq_i}^2 \, d\xx \label{eq:flow-quantization-cost-obj}
%   \end{alignat}
%   \vspace{-\belowdisplayskip}\vspace{-\abovedisplayskip}
%   \begin{empheq}[left={\text{s.t.} \quad \empheqlbrace}]{alignat=1}
%     & \int_{\cell_i} (\xx - \qq_i) \, d\xx = 0 \ \ \forall i , \label{eq:flow-quantization-cost-cons} \\
%     & \cellvol_i = \pvol_i \ \ \forall i . \label{eq:flow-quantization-cost-vol}
%   \end{empheq}
% \end{subequations}
This potential is precisely the joint problem of \S\ref{sec:joint-constraints} with both the sites and the prescribed centroids set to the particles, \ie $\tcentroid_i = \qq_i$.
We can thus express the potential induced by this discretization as
\begin{equation}
  \cpot(\qq) = \max_{\shift, \weight} \jointdual(\shift, \weight) ,
  \label{eq:flow-centroid-potential}
\end{equation}
whose maximum is finite only when the particles $\qq$ are realizable, generating a centroidal power diagram with sites $\qq$ and volumes $\pvol$.

The motion of the particles is then governed by the Lagrangian
\begin{equation}
  \lag(\qq, \dot\qq) = \frac{1}{2} \sum\nolimits_i \pvol_i \norm{\dot\qq_i}^2 - \frac{1}{2\stiff^2}\, \cpot(\qq) ,
  \label{eq:flow-lagrangian}
\end{equation}
where $\stiff \!>\! 0$ sets the strength of the potential.
Since \eqref{eq:flow-centroid-potential} is feasible only when the particles are the centroids of a partition with prescribed volumes, the dynamical system deduced from \eqref{eq:flow-lagrangian} is subject to a \emph{holonomic} constraint, \ie a constraint on the positions alone, independent of the velocities.
% 
The internal forces are computed via the potential gradient relative to $\qq$, which takes the form
\begin{equation}
  \partial_{\qq_i} \cpot =
  % \partial_{\qq_i} \jointdual = 
  2\,\pvol_i\,\shift_i = 2\,\pvol_i\,(\dual_i - \qq_i) ,
  \label{eq:flow-potential-gradient}
\end{equation}
where only the target derivatives contribute, since the site derivatives $\smash{2\,\pvol_i\,(\qq_i - \centroid_i)}$ vanishes when the centroids are the particles.
 
However, the maximizing shifts are only defined up to the homotheties \eqref{eq:joint-nullspace}, so this gradient depends on the choice of generators within the normal cone of $\achievable$ at $\qq$, as noted earlier in \S\ref{sec:background}.
Given the generators $\dual$ at the maximizer, their homotheties are of the form $\smash{s\,\dual_i + \translation}$, with a dilation $s \!>\! 0$ and a translation $\translation \!\in\! \R^d$ as gauge unknowns.
The Euler--Lagrange equations of \eqref{eq:flow-lagrangian} thus read
\begin{equation}
  \ddot\qq_i = \frac{1}{\stiff^2} \bigl( \qq_i - s\,\dual_i - \translation \bigr) ,
  \label{eq:flow-continuous}
\end{equation}
where the gauge $(s, \translation)$ is the multiplier of the holonomic constraint, determined by the requirement that the particles remain realizable.
The potential thus pushes the particles outward, and the constraint absorbs the normal component of this push.
Its tangential component vanishes only when $\qq$ is parallel to $\dual$, \ie when the particles are themselves generators of their power diagram, up to a homothety.
Therefore, the configurations at which the net force in \eqref{eq:flow-continuous} vanishes are the centroidal power diagrams with prescribed volumes, so the potential acts as a regularization of the cells.

\subsection{Constrained Integrator}
\label{sec:flow-constrained}

To integrate \eqref{eq:flow-continuous} in time, we denote the time increment by $\tstep$, and use a superscript to indicate a time step, \eg $\smash{\qq^k}$ refers to the particle positions at time $\smash{k\tstep}$.
Following \citet{marsden2001discrete}, we discretize the action of \eqref{eq:flow-lagrangian} with the trapezoidal rule on the potential, and make it stationary under the holonomic constraint.
The resulting discrete Euler--Lagrange equations simplify to a central finite difference of \eqref{eq:flow-continuous}, taking the form
% , with the velocities $\smash{\vel^k = (\qq^k - \qq^{k-1})/\tstep}$, 
\begin{equation}
  \qq^{k+1}_i - 2\,\qq^k_i + \qq^{k-1}_i = \frac{\tstep^2}{\stiff^2} \bigl( \qq^k_i - s^k\,\dual^k_i - \translation^k \bigr) ,
  \label{eq:flow-discrete}
\end{equation}
where $\smash{\dual^k}$ are generators of the power diagram with centroids $\qq^k$, while the gauge $\smash{(s^k, \translation^k)}$ needs to be chosen so that the new particles $\smash{\qq^{k+1}}$ are realizable.
% Note that the translation $\translation$ are fully determined by the dilation $s$ since summing \eqref{eq:flow-discrete} with weights $\pvol$ must preserve \eqref{eq:partition-feasibility-centroid}.
This update rule corresponds to the Störmer--Verlet scheme in its position form, which is an explicit, second-order, symplectic and time-reversible integrator \cite{Hairer2006GNI}, moving particles in time via
\begin{subequations}
  \label{eq:flow-shake-step}
  \begin{empheq}[left={\empheqlbrace\,}]{align}
    \pred^{k+1} &= \Bigl( 1 + \frac{\tstep^2}{\stiff^2} \Bigr) \qq^k + \tstep\, \vel^k , \label{eq:flow-shake-step-pred} 
    \\[1mm]
    \qq^{k+1} &= \pred^{k+1} - s\,\dual^k - \translation , \label{eq:flow-shake-step-corr}
  \end{empheq}
\end{subequations}
\ie a prediction $\smash{\pred^{k+1}}$ followed by a correction along a homothety of the current generators $\smash{\dual^k}$, where we absorbed the factor $\smash{\tstep^2/\stiff^2}$ into $\smash{(s, \translation)}$,
 % = \tfrac{\tstep^2}{\stiff^2} (s^k, \translation^k)}$
and the new velocities follow as $\smash{\vel^{k+1} \!=\! (\qq^{k+1} \!-\! \qq^k)/\tstep}$.

Instead of a line search over the unknowns $(s, \translation)$ along $\dual^k$ until the new particles $\smash{\qq^{k+1}}$ are realizable, we recast the correction step as a single concave maximization of $\jointdual$ with sites $\qq^k$, prescribed volumes $\pvol$, and prescribed centroids $\smash{\pred^{k+1}}$, while handling the gauge variables via linear constraints on the shifts:
\begin{subequations}
  \label{eq:flow-line-constraint}
  \begin{alignat}{2}
    & \max_{\shift, \weight} \quad && \jointdual(\shift, \weight) \label{eq:flow-line-constraint-obj} \\
    & \text{s.t.} \quad && \sum\nolimits_i \pvol_i\,\shift_i\tr \dual^k_i = 0 , \label{eq:flow-line-constraint-dil} \\
    & && \sum\nolimits_i \pvol_i\,\shift_i = 0 . \label{eq:flow-line-constraint-trans}
  \end{alignat}
\end{subequations}
Introducing multipliers $2s$ and $2\translation$, the stationary condition of \eqref{eq:flow-line-constraint} is exactly the update \eqref{eq:flow-shake-step-corr}, with the new particles as centroids of the resulting power diagram.
Note that, for convenience, we scaled both multipliers by cancel the factor two in \eqref{eq:flow-potential-gradient}.
%
As described in Appendix~\ref{sec:appendix-regularized}, solving \eqref{eq:flow-line-constraint} selects the scale $s$ as the smallest correction that brings the prediction back to $\achievable$.
The constraints \eqref{eq:flow-line-constraint-dil} and \eqref{eq:flow-line-constraint-trans} thus remove the shift components of the dilation and translation modes \eqref{eq:joint-nullspace}, so the maximizer is unique, when it exists, up to the additive constant on the weights.
% \ie the root of the line search closest to the prediction, without any bisection.
%
% Its generators are a homothety of the generators of the new power diagram, so the direction $\smash{\dual^{k+1}}$ of the next step is a by-product of the same solve.

We solve the constrained maximization in \eqref{eq:flow-line-constraint} via a Newton's method, as summarized in Algorithm~\ref{alg:flow-constrained}.
Starting from zero shifts, which satisfy the linear constraints, and weights reset to the latest shifted weights ($\smash{\weight^{k+1} \!=\! \combweight^k}$), each iteration computes the Newton system of $\jointdual$ (\S\ref{sec:joint-constraints}) bordered by the $d + 1$ rows of \eqref{eq:flow-line-constraint-dil} and \eqref{eq:flow-line-constraint-trans}, which returns the updates of the shifts and weights together with the multipliers $(s, \translation)$.
As described in Appendix~\ref{sec:appendix-bordered}, each Newton step costs $d \!+\! 2$ solves of a sparse linear system with the Hessian of $\jointdual$, followed by a dense $(d \!+\! 1) \times (d \!+\! 1)$ system to compute the gauge $(s, \translation)$.
At convergence, the centroids of the maximizer are the new particles $\smash{\qq^{k+1}}$, and its generators $\smash{\qq^k \!+\! \shift}$ provide the direction $\smash{\dual^{k+1}}$ of the next step at no extra cost.

\subsection{Penalty-based Integrator}
\label{sec:flow-penalty}

The constrained integrator in \S\ref{sec:flow-constrained} requires the particles to be realizable at every step, and its correction exists only if the prediction can be brought back to $\achievable$ along a homothety of the current generators.
We now relax this holonomic constraint by representing the particles simply as the \emph{sites} of the power diagram rather than its centroids, while penalizing the centroid constraints in the least-squares sense.
This is precisely the softened problem of \S\ref{sec:regularized-joint} with both the sites and the prescribed centroids at the particles. 
Using $\jointdual_\alpha$ in place of $\jointdual$ in \eqref{eq:flow-centroid-potential}, we thus define the softened potential
\begin{equation}
  \cpot_\alpha(\qq) = \max_{\shift, \weight} \jointdual_\alpha(\shift, \weight) ,
  \label{eq:flow-softened-potential}
\end{equation}
% which replaces the centroid constraint in $\cpot$ with the volume-weighted least-squares residual $\smash{\alpha \sum_i \pvol_i \norm{\centroid_i - \qq_i}^2}$ between the centroids and the particles.
Since $\jointdual_\alpha$ admits a unique maximizer for any targets, the potential in \eqref{eq:flow-softened-potential} is feasible for any distinct particles, so the Lagrangian \eqref{eq:flow-lagrangian} with $\cpot_\alpha$ in place of $\cpot$ is unconstrained.

By \eqref{eq:regularized-projection}, the cells of the maximizer of \eqref{eq:flow-softened-potential} have centroids $\centroid$ at the projection onto $\achievable$ of the dilated particles $(1 + 1/\alpha)\,\qq$.
By \eqref{eq:regularized-gradient-shift}, the optimal shifts are then $\smash{\shift_i = \alpha\,(\qq_i - \centroid_i)}$, which vanish only when $\centroid = \qq$, so the particles are in general no longer the centroids of their cells.
%
The gradient of this potential is
\begin{equation}
  \partial_{\qq_i} \cpot_\alpha
  % = \partial_{\qq_i} \jointdual_\alpha
  % = 2\,\pvol_i\,(\qq_i - \centroid_i) + 2\,\pvol_i\,\shift_i
  = 2\,\Bigl( 1 + \frac{1}{\alpha} \Bigr)\, \pvol_i\,\shift_i ,
  \label{eq:softened-site-gradient}
\end{equation}
which, unlike \eqref{eq:flow-potential-gradient}, has no gauge, since the maximizer of $\jointdual_\alpha$ is unique.
Using it as the force in the derivation of \S\ref{sec:flow-constrained}, with the velocities $\smash{\vel^{k+1} \!=\! (\qq^{k+1} \!-\! \qq^k)/\tstep}$, the discrete Euler--Lagrange equations become a velocity update followed by a position update,
\begin{subequations}
  \label{eq:flow-softened-step}
  \begin{empheq}[left={\empheqlbrace\,}]{align}
    \vel^{k+1}_i &= \vel^k_i + (1 + \alpha)\, \frac{\tstep}{\stiff^2} \bigl( \centroid^k_i - \qq^k_i \bigr) , \label{eq:flow-softened-step-vel}
    \\[1mm]
    \qq^{k+1}_i &= \qq^k_i + \tstep\, \vel^{k+1}_i , \label{eq:flow-softened-step-pos}
  \end{empheq}
\end{subequations}
where $\centroid^k$ are the centroids of the power diagram that maximizes $\jointdual_\alpha$ with sites and target centroids at $\qq^k$, and the force term equals $\smash{-(1 + 1/\alpha)\,(\tstep/\stiff^2)\,\shift^k_i}$ by \eqref{eq:regularized-gradient-shift}.
This update is symplectic and time-reversible as in \S\ref{sec:flow-constrained}, at the cost of each particle drifting from the centroid of its cell by $\smash{\shift_i/\alpha}$.


Each step thus evaluates the force at the current particles by a single unconstrained maximization of $\jointdual_\alpha$, with neither a line search of \eqref{eq:flow-shake-step-corr} nor the linear constraints of \eqref{eq:flow-line-constraint}.
As before, we solve this maximization via a Newton's method, starting from zero shifts and the weights reset to the latest shifted weights.
However, the Newton system of our softened objective needs no border, and its only null-space mode is the additive constant on the weights.
At convergence, the force \eqref{eq:softened-site-gradient} follows directly from the optimal shifts.
Note that the stiffness $\alpha$ is kept fixed over time, so the bound \eqref{eq:regularized-distinct} is only a sufficient condition for the maximizer to be a power diagram.
% TODO: behavior when the generators collide for large alpha; numerical validation in ../src.
We summarize this integrator in Algorithm~\ref{alg:flow-penalty}.

\begin{algorithm}[t]
  \caption{Penalty-based integrator (\S\ref{sec:flow-penalty}).}
  \label{alg:flow-penalty}
  \begin{algorithmic}[1]
    \Require volumes $\pvol$, particles $\qq^0$, velocities $\vel^0$, time step $\tstep$, strength $\stiff$, stiffness $\alpha \!>\! 0$, tolerance $\tau$
    \State $\weight \gets 0$
    \For{$k = 0, 1, 2, \dots$}
      \LineComment{Volume solve with sites $\qq^k$}
      \State $\weight \gets \argmax_\weight \dualenergy(\weight)$, by Newton's method from $\weight$
      \LineComment{Newton solve of $\jointdual_\alpha$ with sites and targets $\qq^k$}
      \State $\shift \gets 0$
      \Repeat
        \State solve the Newton system of $\jointdual_\alpha$ for $(\delta\shift, \delta\weight)$, with one weight pinned
        \State $(\shift, \weight) \gets (\shift, \weight) + \omega\,(\delta\shift, \delta\weight)$, halving $\omega$ from $1$ until no cell is empty
      \Until{$\norm{(\delta\shift, \delta\weight)} < \tau$}
      \LineComment{Update \eqref{eq:flow-softened-step}, with the force \eqref{eq:softened-site-gradient}}
      \State $\vel^{k+1} \gets \vel^k - (1 + 1/\alpha)\,(\tstep/\stiff^2)\, \shift$
      \State $\qq^{k+1} \gets \qq^k + \tstep\,\vel^{k+1}$
    \EndFor
  \end{algorithmic}
\end{algorithm}

\subsection{Discussion}
\label{sec:flow-discussion}

We conclude this section by comparing our two integrators, and by relating them to SHAKE, to the scheme of \citet{Gallouet2018Lagrangian}, and to Power Particles \cite{deGoes2015PowerParticles}.

\paragraph{Constrained vs.\ relaxed integrators.}
Both our integrators are variational and solved by Newton's method on a concave objective, but they use this maximization for different purposes.
The constrained integrator maximizes $\jointdual$ under the linear constraints \eqref{eq:flow-line-constraint} to solve for the homothety gauges $(s, \translation)$ of its correction, which replaces the line search along $\dual^k$, but has no solution when the prediction cannot be brought back to $\achievable$ along this direction.
The penalty-based integrator instead maximizes $\jointdual_\alpha$ only to evaluate the force \eqref{eq:softened-site-gradient} at the current particles, and has a unique maximizer for any distinct particles.
As $\alpha \!\to\! \infty$, the maximizer of $\jointdual_\alpha$ approaches the one of $\jointdual$, but the potential $\cpot_\alpha$ becomes stiff, which restricts the time step of this explicit scheme.
%
Neither integrator constrains the velocities to be divergence-free, instead the weights keep the prescribed volumes exactly at every step and the velocities are updated by the gradient of the potential, which takes the role of the pressure force.
 % (\S\ref{sec:regularized-joint}), the gradient \eqref{eq:softened-site-gradient} recovers \eqref{eq:flow-potential-gradient}, and the update \eqref{eq:flow-softened-step} recovers \eqref{eq:flow-discrete} with the dilation $s^k \!=\! 1$ and the translation $\smash{\translation^k \!=\! 0}$.
% TODO: step-size restriction of the penalty scheme as a function of alpha.

\paragraph{Link to {SHAKE}.}
The constrained integrator of \S\ref{sec:flow-constrained} is an instance of the SHAKE algorithm \cite{Ryckaert1977Shake}, \ie the Störmer--Verlet scheme for a holonomic constraint.
Unlike classical applications of SHAKE, which iterate over explicit constraint functions, our holonomic constraint is a single one whose normal at the particles is given by their generators $\smash{\dual^k}$ up to homotheties.
The constraint force is thus known in direction, and the correction reduces to a single constrained maximization, which finds its magnitude without iterating over constraints.
The same holds for RATTLE \cite{Andersen1983Rattle}, the variant of SHAKE that additionally keeps the velocities tangent to the constraint.
In our settings, this velocity correction reduces to a single projection of $\smash{\vel^{k+1}}$ onto the orthogonal complement of the generators $\smash{\dual^{k+1}}$ and the translations, \ie $\smash{\sum_i \pvol_i\,(\dual^{k+1}_i)\tr \vel^{k+1}_i = 0}$ and $\smash{\sum_i \pvol_i\,\vel^{k+1}_i = 0}$, computed in closed form from the generators already returned by the position correction.

\paragraph{Link to \citet{Gallouet2018Lagrangian}}
Our penalty-based integrator in \S\ref{sec:flow-penalty} generalizes the method of \citet{Gallouet2018Lagrangian}.
For $\alpha \!=\! 0$, the potential \eqref{eq:flow-softened-potential} reduces to b$\smash{\cpot_0(\qq) \!=\! \max_\weight \dualenergy(\weight)}$, which is the potential of \citet{Gallouet2018Lagrangian} with the factor $1/(2\stiff^2)$.
Consequently, the integrator in \cite[Eq.~(1.11)]{Gallouet2018Lagrangian} gives exactly the update of \eqref{eq:flow-softened-step} with $\alpha \!=\! 0$.
% 
As discussed in \citet{Gallouet2018Lagrangian}, this flow convergences towards smooth solutions of the Euler equations for $\tstep \!\le\! \stiff$, with an error bound that vanishes only for $\tstep \!=\! o(\stiff^2)$.
% 
However, their scheme ties particles to the centroids only through the strength $\stiff$ of the potential, so the velocities are divergence-free only in the limit $\stiff \!\to\! 0$.
Instead, our stiffness $\alpha \!>\! 0$ penalizes the centroid residual directly, and thus controls the relaxation of the holonomic constraint independently of $\stiff$.

\paragraph{Link to \citet{deGoes2015PowerParticles}}
The Power Particles method \cite{deGoes2015PowerParticles} advances sites $\qq^k$ and velocities $\vel^k$ via an operator splitting scheme whose steps relate to our penalty-based integrator.
First, they enforce the volumes by maximizing $\dualenergy$ of \eqref{eq:volume-constrained} over the weights for the sites $\qq^k$, which is equivalent to maximizing $\jointdual_\alpha$ with fixed shifts.
Next, they project the velocities $\vel^k$ onto those that preserve the volumes to first order by solving a Poisson equation, which produces the new velocities $\smash{\vel^{k+1}}$ used in place of $\vel^k$ in the update \eqref{eq:flow-softened-step}.
This projection resembles a single Newton step on $\jointdual_\alpha$ from zero shifts and converged weights, with target centroids $\smash{\tcentroid \!=\! \centroid^k \!+\! \tstep\,\vel^k/\alpha}$.
In particular, the discrete divergence and gradient operators derived by \citet{deGoes2015PowerParticles} correspond to the shift--weight block \eqref{eq:joint-hessian} of its Hessian.
Their pressure projection is then reproduced by lumping the shift--shift block to the mass matrix $\smash{\mathrm{diag}(\pvol_i)}$ and zeroing out the weight--weight block.
The Schur complement of the resulting Newton system is the Poisson equation, in which the weight update plays the role of the pressure, while the shift update is the pressure correction that yields $\smash{\vel^{k+1}}$.
Finally, for the particular stiffness $\smash{\alpha = \stiff^2/\tstep^2 - 1}$, the velocity update \eqref{eq:flow-softened-step-vel} adds the Lloyd displacement $\smash{(\centroid^k - \qq^k)/\tstep}$ to $\vel^k$, so the position update \eqref{eq:flow-softened-step-pos} becomes $\smash{\qq^{k+1} = \centroid^k + \tstep\,\vel^k}$, \ie a full Lloyd step onto the centroids followed by advection.
With the pressure-corrected velocities $\smash{\vel^{k+1}}$ in place of $\vel^k$, this is exactly their site update $\smash{\qq^{k+1} \!=\! \centroid^k \!+\! \tstep\,\vel^{k+1}}$, although their stored velocities $\smash{\vel^{k+1}}$ omit the Lloyd displacement.
Since this stiffness is tied to the time step, the splitting does not discretize a Lagrangian, in contrast to our formulations.

\end{document}
